\section{The demos as experiments}
\label{sec:demos}

Each demo below is the acceptance test of one part of the plan. For each we fix, before building it, what we
expect to measure and what result would make us drop or redesign that part; the predictions are numbered so
that the next paper can report against them, as \texttt{experiments/round2/README.md} does for round 2. They
are ordered by what each one depends on, and by cost: the first is days, the last is open.

\subsection{An odometer of digit networks}
\label{sec:demo-odometer}

\emph{Setup.} The digit block and machine of \cref{fig:odometer-spec}, once compiled and once trained; three
digits checked exhaustively over all $1{,}000$ states and both actions; six digits sampled over $10^5$ random
increment sequences of length up to $10^4$ and compared with integer arithmetic.

\begin{prediction}
\label{pred:odometer}
Both odometers are exact: no mismatch with the composed oracle and no violation on any of the $2{,}000$ cells
of the three-digit machine or on any sampled sequence of the six-digit one, and the trained digit's network is
the same $\sim80$k-parameter network for every instance.
\end{prediction}
\emph{Falsifier.} Any carry error. A carry error would mean \cref{thm:compose}'s premise failed in practice
(a block exact on its cells is not exact when the action arrives as an event), and the wiring would need a
different semantics before anything larger is built.

\subsection{A machine that runs itself}
\label{sec:demo-chain}

\emph{Setup.} A clock block drives a chain of finite blocks: a toggle flips on every advance; its flip to
\texttt{on} increments a digit (the odometer's block); the digit's wrap moves an item one place in a reorder
list; the list returning to its initial order turns the three-way switch one place clockwise; the switch
coming back to \texttt{r} resets the digit. Three networks from \cref{tab:round2}, the digit, and one clock,
five wires, rendered as the machine it is, with every hand-off checked live. Continuous blocks join in
\cref{sec:demo-ball}, not here, so that agreement with the composed oracle is exact rather than within a
tolerance.

\begin{prediction}
\label{pred:chain}
The chain runs for $10^4$ ticks with no violation at any block and no disagreement with the composed oracle,
at under $16$\,ms per tick on a laptop CPU.
\end{prediction}
\emph{Falsifier.} A violation anywhere, or a tick that cannot fit the frame. The former says a hand-off delivered
a state a block was not trained for (a gap in the composed spec's reachable set that the block's spec did not
cover); the latter says the budgets of \cref{sec:budgets} are wrong.

\subsection{Grab anything while it runs}
\label{sec:demo-grab}

\emph{Setup.} The same chain; at random ticks a scripted pointer grabs a list item or the switch's knob, drags
it along a recorded path and releases; the machine's state $k$ ticks after release is compared with the
composed oracle's prediction from the dropped state, for $k\in\{1,10,100\}$ over $1{,}000$ grabs.

\begin{prediction}
\label{pred:grab}
Agreement is $100\%$ at every $k$: a drag mid-run is just another external action to the composed spec, and the
blocks downstream re-plan from the dropped state exactly.
\end{prediction}
\emph{Falsifier.} Any disagreement, which would localise to either the detach-and-queue semantics of
\cref{sec:livedrag} or to a drag future the runtime offered that the composed spec does not.

\subsection{The desktop board}
\label{sec:demo-board}

\emph{Setup.} Cards on a grid; a board block, a selection block, a focus block, a menu block and the undo
wrapper, under the bindings of \cref{fig:bindings}. Tests are replayed traces: $1{,}000$ random input sequences
of length $50$ drawn from the binding table (keys with modifiers, clicks, drags, marquees), replayed through the
runtime and through the composed oracle.

\begin{prediction}
\label{pred:board}
Runtime and composed oracle agree on every trace, and the machine answers every input within the frame.
Selection and focus as separate blocks keep the board's network no larger than kanban's.
\end{prediction}
\emph{Falsifier.} Any divergence (a bug in bindings, selection semantics or the menu's guards), or a board
network that only reaches exactness when selection is folded into it, which would say the read-coupling
through item sets is tighter than \cref{sec:couplings} assumes.

\subsection{A ball along track blocks}
\label{sec:demo-ball}

\emph{Setup.} Three continuous track blocks in a chain; a ball's position is continuous within a block; the
hand-off from one track to the next is an event with and without the snapping rule of \cref{sec:handoff}.
Measure the error in position against the composed oracle after $1$, $10$ and $100$ hand-offs over $200$ runs.

\begin{prediction}
\label{pred:ball}
Without snapping, error grows roughly linearly in the number of hand-offs and exceeds $1\%$ of range within
ten; with snapping it is bounded by one block's within-episode error ($\le1\%$) at every count.
\end{prediction}
\emph{Falsifier.} Growth with snapping on, which would mean the error enters inside a block's episode faster
than an episode ends, and the continuous block itself needs the compiled-core treatment of \cref{sec:compiled}
or a shorter episode.

\subsection{Order and cost}
The odometer first, because it tests the theorem with the least machinery and needs only \texttt{emits},
\texttt{machine} and the composed oracle. Then shapes in the render block and the chain, which is the first
thing a viewer can watch; then grabbing, which needs nothing new beyond the clock detach; then the board, which
needs bindings, the selection and focus blocks and the undo wrapper; then the ball, which needs more than one
continuous field per block and the snapping rule. Each is a demo in the lab on the day it passes its test.
